% ============================================================
\section{Periodic Cauchy integrals}\label{app:cauchy}
% ============================================================

It is convenient to work with complex densities. For $z\in\Zc^{1,\gamma}$ and $G\in C^{0,\gamma}(\T;\C)$ set
\[
\mathfrak c[z;G](x):=\frac1{4\pi i}\int_\T G(\beta)\,\partial_\beta z(\beta)\cot\frac{z(\beta)-x}{2}\dd\beta,\qquad x\in\C\setminus\Gamma,
\]
so that $\Cf[z;f]^\flat=\mathfrak c[z;f^\flat]$.

\subsection{Reduction to closed curves}
With $w=e^{ix}$ and $\omega=e^{i\zeta}$, a direct computation gives
\[
\tfrac12\cot\tfrac{\zeta-x}{2}\dd\zeta=\frac{\dd\omega}{\omega-w}-\frac i2\dd\zeta .
\]
The map $x\mapsto e^{ix}$ sends the cylinder $\T\times\R$ to $\C\setminus\{0\}$, $\Gamma$ to a closed Jordan curve $\gamma$ of class $C^{k,\gamma}$ winding once around the origin, $\Omega^+(z)$ to the
punctured interior of $\gamma$ and $\Omega^-(z)$ to its exterior. Hence $\mathfrak c[z;G]$ is the classical Cauchy integral $\frac1{2\pi i}\int_\gamma\frac{\tilde G(\omega)\dd\omega}{\omega-w}$ of the
transported density $\tilde G$, minus the constant $\frac1{4\pi}\int_\T G\partial_\beta z\dd\beta$. The classical Plemelj--Privalov theorem for closed $C^{1,\gamma}$ curves
\cite[Chapter~2]{Mus53}, \cite[Chapter~7]{Kre14} then gives Lemma~\ref{lem:cauchy}(b): $\mathfrak c[z;G]$ extends to a $C^{0,\gamma}$ function on each closed region, with norm bounded by
$C\|G\|_{C^{0,\gamma}}$, where $C$ depends only on $\gamma$, $\delta(z)$ and $\|z\|_{\Zc^{1,\gamma}}$; its boundary values are given by the Plemelj formulas
$\mathfrak c^\pm=\pm\frac12G+\pv\,\mathfrak c$. The value at $w=0$ corresponds to the limit $x_2\to+\infty$.

\subsection{Derivative identities}
\begin{lemma}\label{lem:identities}
Let $z\in\Zc^{2,\gamma}$ and $G\in C^{1,\gamma}(\T;\C)$. Then, in $\C\setminus\Gamma$,
\begin{equation}\label{eq:dx}
\partial_x\,\mathfrak c[z;G]=\mathfrak c\Big[z;\frac{\partial_\beta G}{\partial_\beta z}\Big].
\end{equation}
If $z=z(\beta,t)$ and $G=G(\beta,t)$ are $C^1$ in $t$ with values in $\Zc^{2,\gamma}$ and $C^{1,\gamma}$, then
\begin{equation}\label{eq:dt}
\partial_t\,\mathfrak c[z;G](x)=\mathfrak c\Big[z;\partial_tG-\frac{\partial_\beta G}{\partial_\beta z}\,\partial_tz\Big](x),\qquad x\notin\Gamma(t).
\end{equation}
\end{lemma}

\begin{proof}
Since $\partial_\beta\cot\frac{z(\beta)-x}2=-\frac12\csc^2\frac{z(\beta)-x}2\,\partial_\beta z$, we have $\partial_x\cot\frac{z(\beta)-x}2=-\partial_\beta\cot\frac{z(\beta)-x}2/\partial_\beta z$ and
$\partial_t\cot\frac{z(\beta)-x}2=\partial_\beta\cot\frac{z(\beta)-x}2\cdot\partial_tz/\partial_\beta z$. Integrating by parts in $\beta$ (there are no boundary terms by periodicity),
\[
\partial_x\mathfrak c=-\frac1{4\pi i}\int G\,\partial_\beta\cot\dd\beta=\frac1{4\pi i}\int\partial_\beta G\cot\dd\beta,
\]
which is \eqref{eq:dx}, and
\[
\partial_t\mathfrak c=\frac1{4\pi i}\int\Big[(\partial_tG\,\partial_\beta z+G\,\partial_\beta\partial_tz)\cot-\partial_\beta(G\partial_tz)\cot\Big]\dd\beta
=\frac1{4\pi i}\int\big(\partial_tG\,\partial_\beta z-\partial_\beta G\,\partial_tz\big)\cot\dd\beta,
\]
which is \eqref{eq:dt}.
\end{proof}

Iterating \eqref{eq:dx} and using the Plemelj--Privalov theorem gives Lemma~\ref{lem:cauchy}(d): if $z\in\Zc^{j+1,\gamma}$ and $G\in C^{j,\gamma}$ then all derivatives of order $\le j$ of $\mathfrak c[z;G]$ are
Cauchy integrals of $C^{0,\gamma}$ densities, hence $C^{0,\gamma}$ up to $\Gamma$. If $x=x(t)$ is a $C^1$ path in $\Omega^\pm(z(t))$, then since $\mathfrak c$ is holomorphic,
\begin{equation}\label{eq:chain}
\frac{\dd}{\dd t}\,\mathfrak c[z;G]\big(x(t)\big)=\mathfrak c\Big[z;\partial_tG-\frac{\partial_\beta G}{\partial_\beta z}\partial_tz\Big](x)+\mathfrak c\Big[z;\frac{\partial_\beta G}{\partial_\beta z}\Big](x)\;\dot x,
\end{equation}
and by uniform convergence of the derivatives up to $\Gamma$ the same formula holds for the boundary values $t\mapsto\mathfrak c[z;G]^\pm(\alpha)$ with $x(t)=z(\alpha,t)$.

\subsection{Continuity}
\begin{lemma}\label{lem:continuity}
Let $z_k\to z$ in $\Zc^{1,\gamma}$ and $G_k\to G$ in $C^{0,\gamma}(\T;\C)$. Then $\mathfrak c[z_k;G_k]^\pm\to\mathfrak c[z;G]^\pm$ uniformly on $\T$, and if $x_k\in\overline{\Omega^+(z_k)}$ with $x_k\to z(\alpha)$, then
$\mathfrak c[z_k;G_k](x_k)\to\mathfrak c[z;G]^+(\alpha)$; similarly from below.
\end{lemma}

\begin{proof}
By the Plemelj--Privalov theorem the functions $\mathfrak c_k:=\mathfrak c[z_k;G_k]$ are uniformly $C^{0,\gamma}$ on $\overline{\Omega^+(z_k)}$, with a constant $L$ independent of $k$. Given $\eps>0$ let
$y=z(\alpha)+\eps n(\alpha)$; for $k$ large $y\in\Omega^+(z_k)$ at distance $\ge\eps/2$ from $\Gamma_k$, so $\mathfrak c_k(y)\to\mathfrak c[z;G](y)$ by dominated convergence. Hence
\[
\limsup_k|\mathfrak c_k(x_k)-\mathfrak c[z;G]^+(\alpha)|\le\limsup_k\big(L|x_k-y|^\gamma+L\eps^\gamma\big)\le 2L\eps^\gamma,
\]
and $\eps$ is arbitrary. Uniformity in $\alpha$ follows from the same argument and the equicontinuity.
\end{proof}

\subsection{Proof of Lemma~\ref{lem:timereg} and of the regularity of $G$ in Lemma~\ref{lem:accel}}
By Lemma~\ref{lem:geometry}, the curves $z^\pm(\cdot,t)$ are polynomials in $t$ with coefficients in $C^{3,\gamma}$, with $\delta(z^\pm)\ge\frac12\delta(z_0)$ on $[0,T_1]$. The normals $n^\pm$, the lengths
$|\partial_\alpha z^\pm|$ and the velocities $V^\pm=\partial_tz^\pm\cdot n^\pm$ are therefore smooth in $t$ with values in $C^{2,\gamma}$, and so are the complex densities $G^\pm:=(f^\pm)^\flat$ and their time
derivatives.

By \eqref{eq:vS}, $(v_S^+)^\flat=\mathfrak c[z^+;G^+]^-+\mathfrak c[z^-;G^-]\circ z^+$, where $z^+(\alpha,t)\in\Omega^+(z^-(\cdot,t))$ for $t>0$, and similarly for $v_S^-$. Applying \eqref{eq:chain} once, $\partial_tv_S^\pm$
is a sum of boundary values and evaluations on the other curve of Cauchy integrals with densities
\[
G_1^\pm=\partial_tG^\pm-\frac{\partial_\beta G^\pm}{\partial_\beta z^\pm}\partial_tz^\pm,\qquad G_2^\pm=\frac{\partial_\beta G^\pm}{\partial_\beta z^\pm},
\]
multiplied by $\partial_tz^\pm$; these densities are bounded in $C^{1,\gamma}$ uniformly in $t$. Applying \eqref{eq:chain} once more, $\partial_t^2v_S^\pm$ involves Cauchy integrals of densities bounded in
$C^{0,\gamma}$, evaluated on $\Gamma^\pm(t)$ from the side of the strip, and these are bounded uniformly in $t\in(0,T_1]$ by the Plemelj--Privalov theorem. Since $F^\pm$ are polynomials in
$v_S^\pm\cdot n^\pm$, $V^\pm$, $n_2^\pm$ and $|\partial_\alpha z^\pm|$, we obtain that $\partial_tF^\pm$ and $\partial_t^2F^\pm$ are bounded in $C^0(\T)$ on $(0,T_1]$. By Lemma~\ref{lem:continuity} applied to the
densities appearing in $F^\pm$ and $\partial_tF^\pm$, these functions extend continuously to $t=0$, which proves the first part of Lemma~\ref{lem:timereg}.

For \eqref{eq:v-strip}, let $x=X(\alpha,\lambda,t)$, so that $|x-z^\pm(\alpha,t)|\le2t\kappa(\alpha)$. The uniform Hölder bounds on $\overline{\Omega^-(z^+)}$ and $\overline{\Omega^+(z^-)}$ give
$|v(x,t)-v_S^+(\alpha,t)|\le Ct^\gamma$, and the Lipschitz bound in time just proved gives $|v_S^+(\alpha,t)-v_S^{(0)}(\alpha)|\le Ct$.

Finally, $G=\partial_tM(\cdot,0)$ for $a=0$, $h=0$ is obtained by letting $t\to0$ in the formula for $\partial_tM$. In the limit all Cauchy integrals become boundary values on $z_0\in\Zc^{6,\gamma}$, and
the densities involve $n_0\in C^{5,\gamma}$, $\partial_tz^\pm(\cdot,0)=(U_0\pm c^\pm)n_0\in C^{4,\gamma}$, their first $\alpha$-derivatives, and $\partial_t^2z^\pm(\cdot,0)=0$. Thus
$G^\pm(\cdot,0)\in C^{4,\gamma}$, $G_1^\pm(\cdot,0),G_2^\pm(\cdot,0)\in C^{3,\gamma}$, and by Lemma~\ref{lem:cauchy}(d) the boundary values are in $C^{3,\gamma}$. The remaining factors ($\partial_tz^\pm(\cdot,0)$, $n_0$,
and the time derivatives at $t=0$ of $n^\pm$, $V^\pm$ and $|\partial_\alpha z^\pm|$, which involve $\partial_\alpha\partial_tz^\pm(\cdot,0)$) belong to $C^{3,\gamma}$. Hence $G\in C^{3,\gamma}(\T)$.
\qed
